\subsection{Interpreting Quality and Alignment}
Text adaptation can change model rankings even with a frozen generation
backbone, as the Motion-X FID reversal illustrates. Source-encoder quality
and adapted-encoder retrieval therefore characterize distinct systems; their
advantages cannot be combined into a joint performance claim.
Within either configuration, FID$_D$ measures a marginal distribution gap,
not pairwise invariance. A smaller gap can reflect changes in either
prompt-conditioned distribution and should be interpreted with reference-based
FID, text alignment, and paired measures from the same generation set.

\subsection{What the Theory Explains}
The local analysis separates a loss-value identity from an optimization claim.
Cosine consistency measures projected Jacobian energy of the auxiliary
representation. Transferring that consistency to the generator requires an
additional representation bridge or a nonzero update aligned with the chosen
risk gradient. The detached implementation lacks the latter gradient path.
Enabling it in a connectivity test establishes reachability, not a beneficial
update or a trained-model result.

The observed reductions in hidden-state response and Jacobian norms are
consistent with lower local sensitivity, but do not establish the directional
ordering in~\eqref{eq:theory_contraction}. The sequence KL bound additionally
requires control at shared prefixes throughout the rollout distribution.
Selected teacher-forced prefixes cannot verify that premise. The theory thus
provides testable conditions for a mechanism, rather than a causal explanation
of every reported gain or a guarantee about FID and sampled trajectories.

\subsection{Limitations and Next Steps}
Three limitations determine the scope of the evidence. First, the generation
tests use one seed; the 20 retrieval partitions quantify candidate grouping,
not generation or training variability. The column-wise source/adapted
presentation was adopted after observing results, and both full configurations
are retained. Evaluator adaptation, lowercase lookup, and architecture-specific
unconditional inputs further limit comparisons with other protocols.

Second, source training budgets, selected endpoints, and active losses differ.
Equal text-adaptation budgets do not remove these confounds. A matched
attribution study should fix initialization, data order, budget, validation
selection, and decoding, then evaluate generator-gradient alignment and
virtual updates across batches using an independently defined risk and the
actual optimizer direction.

Third, the benchmark evaluates preserving rewrites; altering pairs lack edited
motion ground truth and support only internal diagnostics. The next evaluation
should test compliance at the edited factor and preservation elsewhere,
alongside multiple-seed and matched-budget studies. A single consistency
score cannot capture both desired invariance and desired edit sensitivity.